\subsection{Schmidt canonical form on an open chain}
\label{sec:obc_schmidt_canonical}

\begin{lemma}[Weighted Schmidt bases at physical cuts]
    \label{lem:obc_weighted_schmidt_cut_bases}
    \lean{MPSPreparation.exists_weightedSchmidtCutBasis,
        MPSPreparation.exists_canonicalCutBasisFamily}
    \notready
    \uses{def:varying_bond_obc_chain}
    Let $M_k$ be the physical coefficient matrix of a finite state at cut
    $k$, with rank $r_k$. Its compact singular-value decomposition
    $M_k=V_k^\dagger D_kU_k$ gives a basis of its column space consisting
    of the columns of $F_k=V_k^\dagger D_k$. In this basis the suffix
    coefficient matrix is $G_k=U_k$, so
    \[
        F_k^\dagger F_k=D_k^2>0,\qquad G_kG_k^\dagger=I_{r_k}.
    \]
    For a normalized state on a positive number of sites, these bases
    can be chosen simultaneously with endpoint vectors $1$ and $\psi$.
    The endpoint Gram matrices are both the scalar identity, and the
    endpoint suffix factors are row coisometries as well.
\end{lemma}
\begin{proof}
    The inverse $D_k^{-1}V_k$ proves independence of the weighted left
    singular vectors. The equation $F_k=M_kU_k^\dagger$ places them
    in the physical column space, whose dimension is $r_k$. Their
    coordinates reconstruct the right factor $U_k$. At the two
    one-dimensional endpoints choose $1$ and the actual normalized
    state; normalization gives the same unit Gram and suffix equations.
\end{proof}

\begin{lemma}[Physical configuration bijections and endpoint Grams]
    \label{lem:obc_canonical_configurations_endpoints}
    \lean{MPSPreparation.extendCutPrefixEquiv,
        MPSPreparation.extendCutSuffixEquiv,
        MPSPreparation.initialCutSuffixEquiv,
        MPSPreparation.finalCutPrefixEquiv,
        MPSPreparation.cutPrefixMatrix_zero_gram,
        MPSPreparation.cutPrefixMatrix_last_gram,
        MPSPreparation.cutSuffixMatrix_zero_coisometry,
        MPSPreparation.cutSuffixMatrix_last_coisometry}
    \notready
    \uses{lem:obc_weighted_schmidt_cut_bases}
    Adjoining one physical letter gives bijections from the preceding
    prefix configurations times the letter to the next prefix
    configurations, and from the letter times the next suffix
    configurations to the preceding suffix configurations. The final
    prefixes and initial suffixes are each in bijection with the full
    physical configurations. Thus all sums in the endpoint normalization
    and adjacent-cut Gram identities enumerate the actual configurations
    exactly once.
\end{lemma}
\begin{proof}
    Restrict a configuration away from the cut site and record the
    letter at that site. These operations are inverse to adjoining
    that letter. At the endpoints the empty configuration spaces are
    singletons, and the remaining configuration is the full chain.
\end{proof}

\begin{lemma}[Canonical equations from adjacent physical cut factors]
    \label{lem:obc_canonical_cut_recurrences}
    \lean{MPSPreparation.cutBondDensity,
        MPSPreparation.cutBondDensity_step,
        MPSPreparation.cutSuffixMatrix_gram_step,
        MPSPreparation.sum_cutSiteMatrix_mul_conjTranspose_eq_one,
        MPSPreparation.trace_cutBondDensity_step}
    \notready
    \uses{lem:obc_weighted_schmidt_cut_bases,
        lem:obc_canonical_configurations_endpoints}
    Let $A_k^i$ be the coordinate matrices between the chosen adjacent
    cut bases, and put $\Lambda_k=F_k^\dagger F_k$. The actual prefix
    and suffix recursions give
    \[
        \sum_i(A_k^i)^\dagger\Lambda_kA_k^i=\Lambda_{k+1},
        \qquad
        G_kG_k^\dagger=
        \sum_i A_k^i(G_{k+1}G_{k+1}^\dagger)(A_k^i)^\dagger.
    \]
    In particular, if both suffix factors are row coisometries, then
    $\sum_i A_k^i(A_k^i)^\dagger=I$ and
    $\operatorname{tr}\Lambda_{k+1}=\operatorname{tr}\Lambda_k$.
\end{lemma}
\begin{proof}
    Expand the physical prefix and suffix factors by the same coordinate
    slicing matrices, and use the preceding bijections to take their
    Grams. Cyclicity of the trace and the right-canonical equation
    give trace preservation.
\end{proof}

\begin{theorem}[Simultaneous OBC Schmidt canonical form]
    \label{thm:obc_schmidt_canonical}
    \lean{OBCChainTensor.exists_canonical_coeff_eq}
    \notready
    \uses{lem:obc_canonical_cut_recurrences, thm:seqanc_cut_rank_bounds}
    Let $N>0$, and let $\psi$ be a normalized state on $N$ sites of
    physical dimension $d$. There is an exact OBC representation
    $c_A=\psi$ with bonds
    \[
        D_k=\operatorname{rank}_k\psi\le d^{\lfloor N/2\rfloor}
    \]
    and matrices $\Lambda_k\in\mathrm M_{D_k}(\mathbb C)$ satisfying,
    for every site $k$,
    \[
        \sum_i A_k^i(A_k^i)^\dagger=I_{D_k},\qquad
        \sum_i(A_k^i)^\dagger\Lambda_kA_k^i=\Lambda_{k+1}.
    \]
    Both endpoint densities are the scalar identity. Every
    $\Lambda_k$ is diagonal with strictly positive real diagonal
    entries, is positive definite, and has trace one. This is the
    normalized-state OBC canonical form of
    \cite[Theorem~1, conditions~1--3]{PerezGarcia2007Matrix}.
\end{theorem}
\begin{proof}
    Use the weighted Schmidt cut bases from the first lemma in the
    existing exact minimal-bond construction. The endpoint basis vectors
    give the original coefficient tensor without changing its phase.
    The preceding recurrences prove both canonical equations. The Gram
    diagonals are the squares of the retained positive singular values.
    The initial trace is one and is preserved at every site. The physical
    cut-rank bound gives the stated common bond bound.
\end{proof}
